\chapter{Activity specialization}
\label{ch:projects}

\section{Projects as growing networks}

The duplication rule of the previous chapter was born in biology, but copying is
not a biological privilege. This chapter applies the same idea to a very
different system --- the schedule of a large engineering project --- and gets, in
return, a surprising practical payoff: an explanation of why big projects so
reliably run late.

A project schedule can be drawn as a network, called an \emph{activity network}.
Its nodes are the individual activities, or tasks, that make up the project, and
its links record \emph{precedence}: a link from activity $a$ to activity $b$
means $b$ cannot begin until $a$ is finished --- you cannot paint a wall before it
is built. Because work only flows forward in time, the network has no loops; in
the jargon it is a \emph{directed acyclic graph}, with starting activities (no
predecessors) and finishing activities (no successors). For such a network there
is one object that has ruled project management since the 1950s: the
\emph{critical path}. This is the longest chain of dependent activities running
from start to finish --- the route through the project with no slack anywhere
along it, so that delaying \emph{any} activity on it pushes back the whole
project's end date by the same amount \citep{kelley1959}. The manager's
time-honored prescription is therefore simple: watch the critical path, and
guard it at all costs.

And yet projects keep finishing late. A survey of more than ten thousand of them
found only a few percent delivered on time and on budget, and fifteen years of
improvements in management practice barely moved that figure
\citep{flyvbjerg2003}. This stubborn record hints that the critical path --- a
single one-dimensional thread --- may be the wrong thing to watch in schedules
that have grown into genuinely complex networks. Following
\citet{vazquezprojects2023}, this chapter grows such schedules from a local rule,
shows that the critical path \emph{shrinks toward irrelevance} as projects get
bigger, and relocates the real source of delay risk in the network as a whole.

\section{The duplication--split rule}

Where does a schedule come from? Not in one piece. A plan is \emph{refined}: a
coarse activity is broken into finer ones, again and again, until the plan is
detailed enough to execute. There are two natural ways to break an activity in
two, and they correspond to two kinds of work. This is our local rule, with a
single dial, the duplication probability $q$. Beginning from two activities in
series, repeat: pick an existing activity at random and, with probability $q$,
\emph{duplicate} it, otherwise \emph{split} it.
\begin{itemize}
  \item \emph{Duplication} (probability $q$). A \emph{generic} activity is split
    into two smaller activities that run \emph{in parallel} --- think of doubling
    a crew so two teams do the same kind of work at once. Both copies inherit
    \emph{all} the predecessors and \emph{all} the successors of the parent.
    Duplication \emph{widens} the schedule, adding parallel strands of work.
  \item \emph{Split} (probability $1-q$). A \emph{specialized} activity is cut
    into two activities run \emph{in sequence} --- one specialist does the first
    part and hands off to another who does the second. The parent keeps its
    predecessors, the new activity takes over its successors, and a link joins
    the two. Splitting \emph{lengthens} the schedule, extending a chain.
\end{itemize}
\begin{figure}[t]
\centering
\begin{adjustbox}{max width=\linewidth}\begin{tikzpicture}[>={Stealth[round,length=2.2mm]},
  nd/.style={circle,draw=black!65,fill=black!4,minimum size=5.5mm,inner sep=0pt,font=\scriptsize},
  act/.style={circle,draw=accent,fill=black!30,minimum size=5.5mm,inner sep=0pt,font=\scriptsize},
  newl/.style={->,draw=accent,line width=1.2pt}]
  % panel (a) duplication: parallel
  \begin{scope}
    \node[nd]  (p)  at (0,0)      {};
    \node[act] (A1) at (1.4,0.7)  {};
    \node[act] (A2) at (1.4,-0.7) {};
    \node[nd]  (s)  at (2.8,0)     {};
    \draw[newl] (p)--(A1); \draw[newl] (p)--(A2);
    \draw[newl] (A1)--(s); \draw[newl] (A2)--(s);
    \node[font=\footnotesize] at (1.4,-1.7) {(a) duplication: in parallel};
  \end{scope}
  % panel (b) split: series
  \begin{scope}[xshift=5.6cm]
    \node[nd]  (p2) at (0,0)    {};
    \node[act] (B1) at (1.1,0)  {};
    \node[act] (B2) at (2.3,0)  {};
    \node[nd]  (s2) at (3.4,0)   {};
    \draw[newl] (p2)--(B1); \draw[newl] (B1)--(B2); \draw[newl] (B2)--(s2);
    \node[font=\footnotesize] at (1.7,-1.7) {(b) split: in sequence};
  \end{scope}
\end{tikzpicture}\end{adjustbox}
\caption{\textbf{Refining an activity.} A generic activity is replaced either by
(a) two copies that run \emph{in parallel}, both inheriting the predecessor (left,
gray) and successor (right, gray) of the parent --- duplication, which widens the
schedule --- or by (b) two specialized activities that run \emph{in sequence} --- a
split, which lengthens it.}
\label{fig:proj-dupsplit}
\end{figure}

So the dial $q$ measures how a project is organized (Figure~\ref{fig:proj-dupsplit}). Small $q$ means almost every
refinement is a split: work is broken into ever-longer chains of specialists, and
the schedule is nearly a single long line. Large $q$ means almost every
refinement is a duplication: generic work is parallelized into many simultaneous
strands, and the schedule fans out into a broad web. The same copying that
duplicated a gene in the previous chapter here duplicates an activity; the new
ingredient is the \emph{split} --- the act of \emph{specialization} that gives the
chapter its name --- and the forward-flowing, start-to-finish structure of a
schedule.

\section{Predecessors and successors: identical power laws}

Because the rule copies an activity together with its predecessors and
successors, it is a duplication rule of the kind that produced power-law degrees
in the previous chapter --- only now there are two degrees to track, the number of
predecessors (the in-degree) and the number of successors (the out-degree). The
box works out the distribution of predecessors by following how many activities
have each number of them; the same argument applies verbatim to successors. The
result is in the paragraph after.

\begin{calculation}{the distribution of predecessors}
Let $n_k(n)$ be the number of activities with $k$ predecessors when the schedule
has $n$ activities in all. Adding one activity changes this count by
\begin{equation}
\adjustbox{max width=\linewidth}{$\displaystyle
  n_k(n+1) = n_k(n)
   + q\left[\frac{k-1}{n}\,n_{k-1}(n) - \frac{k}{n}\,n_k(n)
            + \frac{1}{n}\,n_k(n)\right]
   + (1-q)\,\delta_{k,1}
$}
  \label{eq:proj-rate}
\end{equation}
the bracket collecting the duplication events (an activity climbs from $k-1$ to
$k$ predecessors when one of its predecessors is duplicated, and so on) and the
last term adding a fresh single-predecessor activity at each split. Looking for a
stable shape $n_k(n)=n\,p_k$ turns this into a recurrence,
\begin{equation}
  p_k = q\big[(k-1)\,p_{k-1} - k\,p_k + p_k\big] + (1-q)\,\delta_{k,1},
  \label{eq:proj-recurrence}
\end{equation}
which solves in closed form (with $\Gamma$ the gamma function) to
\begin{equation}
  p_k = \prod_{s=1}^{k-1}\frac{s}{\tfrac{1}{q}+s}\; p_1
      = \Gamma\!\left(\tfrac{1}{q}+1\right)
        \frac{\Gamma(k)}{\Gamma\!\left(\tfrac{1}{q}+k\right)}\, p_1,
  \label{eq:proj-pk}
\end{equation}
with the power-law tail
\begin{equation}
  p_k \sim k^{-1/q}.
  \label{eq:proj-asym}
\end{equation}
\end{calculation}

The result, Eq.~\eqref{eq:proj-asym}, is a power law with exponent $\gamma=1/q$.
A wide, duplication-heavy schedule (large $q$) has the heaviest tail of
highly-connected activities; at $q=1/2$ we recover the familiar value $\gamma=2$.
The very same reasoning, run on successors instead of predecessors, gives the
\emph{identical} distribution for the out-degree. That coincidence is the model's
signature prediction --- the spread of predecessor counts and the spread of
successor counts should look the same --- and it follows directly from the
even-handed way duplication treats the two ends of an activity.

\section{The vanishing critical path}

The degrees are a local property; the critical path is a global one. Remarkably,
the same rule fixes how the critical path grows with the size of the project, and
the answer overturns the manager's intuition. The box does the short
calculation; the consequences follow.

\begin{calculation}{how the critical path scales}
Let $c$ be the number of activities on the critical path of an $n$-activity
schedule. The path can only lengthen when a \emph{split} lands on it (a split
inserts an activity in series), and the chance that the activity picked for
refinement lies on the critical path is $c/n$. Hence
\begin{equation}
  \frac{dc}{dn} = (1-q)\,\frac{c}{n},
  \label{eq:proj-cp}
\end{equation}
which integrates to $c\sim n^{1-q}$, so the \emph{fraction} of activities on the
critical path shrinks as
\begin{equation}
  \frac{c}{n} \sim n^{-\alpha(q)}, \qquad \alpha(q) \le q.
  \label{eq:proj-cpfrac}
\end{equation}
Because distances here grow as a power of size, $c\sim n^{1/D}$, the network is a
\emph{fractal} with dimension
\begin{equation}
  D = \frac{1}{1-\alpha(q)}.
  \label{eq:proj-fractal}
\end{equation}
\end{calculation}

The headline is Eq.~\eqref{eq:proj-cpfrac}: for any amount of duplication
($q>0$), the critical path is a \emph{vanishing fraction} of the project as the
project grows. The bigger and more parallel the schedule, the smaller the slice
of the work that the manager's cherished path actually covers --- so watching the
critical path means watching an ever-shrinking sliver of the whole. (Simulations
confirm the power-law shrinkage, with the measured exponent $\alpha$ a little
below the simple bound $q$; for instance $\alpha\approx 0.11,\,0.28,\,0.43$ at
$q=0.20,\,0.50,\,0.70$.)

The scaling also tells us what \emph{kind} of network a schedule is. In a
small-world network, distances grow only logarithmically with size --- doubling
the network barely lengthens the longest path. Here distances grow as a
\emph{power} of size, the mark of a \emph{fractal} network: one whose ``size''
and ``length'' are tied together by a fixed dimension $D$, much as a physical
fractal looks the same at every magnification. Duplication--split schedules are
thus a class apart --- power-law degrees like a scale-free network, but power-law
(not logarithmic) distances, unlike a small world. Specialization builds the long
fractal chains; duplication stacks them up in parallel.

\section{Real project schedules}

This is not just a story about a model. \citet{vazquezprojects2023} examined
$77$ real construction projects, together worth more than ten billion dollars,
which yielded $323$ activity networks --- snapshots of projects at different
stages --- ranging from about a hundred to sixteen thousand activities. Three
predictions survive the encounter with data. First, the predecessor and successor
counts are each power-law distributed, and fitting them \emph{separately} yields
duplication indices $q$ that agree, just as the model demands they should.
Second, those indices fall between roughly $0.1$ and $0.5$, most above $0.2$:
real schedules sit on the split-dominated, nearly-linear side --- but with enough
duplication to matter. Third, the share of activities on the critical path falls
with project size right across the collection, confirming that the critical path
really does dwindle as projects grow.

\section{Delay risk as percolation}

If the critical path fades to insignificance and yet projects are still late,
where does the delay risk live? The answer is that it is spread across the whole
network, and it is governed by \emph{percolation} --- the same mathematics that
describes how water seeps through porous rock or how a forest fire spreads from
tree to tree.

Picture an outside shock --- bad weather, a strike, a pandemic --- that pushes
some activities past their planned finish. A delay only passes to a downstream
activity if it is large enough to use up the spare time, the \emph{free float},
built into that dependency; otherwise the buffer absorbs it and the delay stops
there. Now treat each dependency as a channel that passes a delay along with some
probability $p$, and ask what fraction of activities end up delayed beyond their
own slack --- the \emph{delay incidence}. This is exactly a percolation problem:
delays spread through the network of dependencies like fire through a forest, and
percolation theory says such spreading has a sharp \emph{threshold}
\citep{albert2000,ellinas2019,schwartz2002},
\begin{equation}
  p_c = \frac{1}{\langle k\rangle},
  \label{eq:proj-pc}
\end{equation}
where $\langle k\rangle$ is the average number of downstream dependencies per
activity. The meaning of a threshold is stark: below $p_c$ delays stay local and
project-wide overruns essentially never happen; above $p_c$ a finite fraction of
projects suffer a system-wide cascade --- the fire jumps the firebreaks.

The data choose between two stories cleanly. If the critical path governed
delay, the incidence ought to track $p\,c$ (transmission times critical-path
size); in the real projects that correlation is negligible. If percolation
governs delay, the incidence ought to switch on exactly at $p-p_c=0$; and that is
what the construction projects show --- below threshold, delay incidence stays
under one percent; above it, the incidence climbs, reaching fifteen percent of
all activities in the worst cases. The true control knob of project delay is
$p-p_c$, a property of the \emph{whole} network, not $p\,c$, a property of one
path through it.

Here is the chapter's practical lesson, and its link to the rest of the book. The
very same local rule that sets the degree distributions also sets the average
degree $\langle k\rangle$, and through Eq.~\eqref{eq:proj-pc} the threshold for
catastrophic delay. A project's tendency to run late is written into the network
its activities have grown into --- not into any single path through it.

\section{Summary}

Project schedules are forward-flowing networks of activities, grown by a local
refinement rule: pick an activity and either duplicate it into parallel copies
(probability $q$) or split it into specialized sequential parts. This
duplication--split rule predicts identical power-law distributions for the
numbers of predecessors and successors, $p_k\sim k^{-1/q}$; a critical path that
grows only as $c\sim n^{1-q}$ and so vanishes as a fraction of large projects;
and a fractal, rather than small-world, geometry. All three hold in $77$ real
construction projects. And because the critical path becomes irrelevant at scale,
delay risk is better read through percolation, with a threshold $p_c=1/\langle
k\rangle$ fixed by the grown network's average degree --- so that project
performance, in the end, is a property of the activity network as a whole.
